\chapter{Implementation of the Digitisation Pipeline}

The architectural design presented in Chapter~3 is realised through a layered implementation spanning image processing scripts, a deep learning inference pipeline, and an adaptive binarisation module. Each stage is encapsulated in a dedicated Python module with a clearly defined function interface, making the overall system modular and testable. This chapter walks through the implementation of each Phase~1 pipeline stage, covering key functions, configuration parameters, and integration mechanisms, before addressing the ECE signal processing modules and the testing strategy that complete the Phase~1 system.

\section{Technology Stack and Implementation Phases}

\subsection{Software Dependencies}

Table~\ref{tab:tech_stack} summarises the Phase~1 software stack. All library versions are pinned in \texttt{requirements.txt} and \texttt{environment.yml} to ensure reproducible builds.

\begin{table}[htb]
\fontsize{10}{12}\selectfont
\centering
\caption{Technology stack: layer, library, and version}
\label{tab:tech_stack}
\begin{tabular}{|p{3.8cm}|p{4.8cm}|p{3.2cm}|}
\hline
\textbf{Layer} & \textbf{Library} & \textbf{Version} \\
\hline
Image processing & OpenCV & 4.9.0.80 \\
\hline
Image I/O \& EXIF & Pillow & 10.3.0 \\
\hline
Scientific image ops & scikit-image & 0.23.2 \\
\hline
Numerical computing & NumPy & 1.26.4 \\
\hline
AI super-resolution & Real-ESRGAN (basicsr) & 0.3.0 / 1.4.2 \\
\hline
Deep learning runtime & PyTorch & $\geq$ 2.0.0 \\
\hline
\end{tabular}
\end{table}

\subsection{Implementation Phases and Timeline}

The implementation followed a phased schedule aligned with the three-stage pipeline architecture. Table~\ref{tab:impl_phases} summarises each phase, its primary deliverable, estimated duration, and current status.

\begin{table}[htb]
\fontsize{10}{12}\selectfont
\centering
\caption{Implementation phases and completion status}
\label{tab:impl_phases}
\begin{tabular}{|c|p{5.5cm}|p{2.2cm}|p{2.0cm}|}
\hline
\textbf{Phase} & \textbf{Deliverable} & \textbf{Duration} & \textbf{Status} \\
\hline
1 & Environment setup, folder structure, 5 sample images tested & Week 1 & Complete \\
\hline
2 & \texttt{preprocess.py} + \texttt{enhance.py} --- full enhancement pipeline with batch processing & Weeks 2--3 & Complete \\
\hline
3 & \texttt{binarise.py} --- binarisation for all five artefact types & Weeks 4--6 & Complete \\
\hline
\end{tabular}
\end{table}

All three phases are complete and integrated. Further stages are deferred to Phase~2 and beyond.

\section{Stage 1 --- Preprocessing (\texttt{src/preprocess.py})}

Five functions implement the preprocessing chain. They are called in sequence by the main entry point \texttt{preprocess(img\_path, output\_path)}.

\textbf{\texttt{load\_image(path)}:} Reads the image from disk as a BGR NumPy array using OpenCV, then applies \texttt{PIL.ImageOps.exif\_transpose} from the Pillow library to bake the EXIF orientation tag into the pixel data. This step is critical for field photographs taken at arbitrary camera angles: without it, subsequent stages would receive rotated images that would degrade CLAHE tiling, binarisation layout detection, and visual comparison results.

\textbf{\texttt{normalise\_brightness(img)}:} Converts the BGR image to LAB colour space, isolates the L (lightness) channel, and applies OpenCV's CLAHE (\texttt{cv2.createCLAHE(clipLimit=2.0, tileGridSize=(8, 8))}). The enhanced L channel is merged back into LAB and converted to BGR. Operating exclusively on the lightness channel ensures that colour information, which is used by DStretch in Stage~2, is not altered by the brightness normalisation step.

\textbf{\texttt{auto\_white\_balance(img)}:} Implements the grey-world assumption. For each colour channel $c \in \{B, G, R\}$, the channel mean $\mu_c$ is computed. Each channel is multiplied by $\bar{\mu} / \mu_c$, where $\bar{\mu}$ is the mean of all three channel means. The result is clipped to $[0, 255]$. This automated correction removes systematic colour casts without requiring a physical reference target.

\textbf{\texttt{crop\_borders(img, threshold=10)}:} Computes the per-row and per-column pixel mean. Rows and columns whose mean falls below the threshold value of 10 (near-black, corresponding to scanner bed or dark margin) are identified as border regions. The function identifies the innermost non-border rows and columns on all four sides and returns the cropped sub-image.

\textbf{\texttt{process\_directory(input\_dir, output\_dir)}:} Iterates over all supported image files in \texttt{input\_dir} (JPEG, PNG, and TIFF), calls \texttt{preprocess()} for each, and saves the normalised output to \texttt{output\_dir} as JPEG with quality=95.

\section{Stage 2 --- AI Enhancement (\texttt{src/enhance.py})}

\subsection{Denoising}

\textbf{\texttt{denoise(img, strength=10)}:} Wraps OpenCV's \texttt{cv2.fastNlMeansDenoisingColored}, which applies non-local means filtering across all three colour channels. The \texttt{h} parameter (filter strength) is set to the \texttt{strength} argument. ECE noise analysis informed three default values: 10 for stone inscriptions, 15 for palm leaf manuscripts, and 8 for copper plates. Higher values suppress more noise but risk blurring fine character strokes, so the parameter is conservative for artefact types with dense script.

\subsection{Real-ESRGAN Super-Resolution}

\textbf{\texttt{enhance\_with\_realesrgan(img, scale=2, model\_path)}:} Constructs an \texttt{RRDBNet} model with parameters \texttt{num\_in\_ch=3}, \texttt{num\_out\_ch=3}, \texttt{num\_feat=64}, \texttt{num\_block=23}, \texttt{num\_grow\_ch=32} matching the pre-trained \texttt{RealESRGAN\_x4plus.pth} architecture \cite{Wang2021}. The \texttt{RealESRGANer} inference wrapper is configured with:
\begin{itemize}
    \item \texttt{tile=400} --- tiles large images into 400-pixel blocks before inference to prevent GPU out-of-memory errors on high-resolution inscription scans.
    \item \texttt{tile\_pad=10} --- adds a 10-pixel overlap between tiles to avoid visible seam artefacts at tile boundaries.
    \item \texttt{half=False} --- uses full 32-bit floating point precision, avoiding numerical instability on inscriptions with very low contrast gradients.
    \item \texttt{outscale=2} --- outputs at 2$\times$ magnification from the internally 4$\times$ model, preventing hallucinated detail on dense ancient character sets.
\end{itemize}

For line-art carving styles, \texttt{RealESRGAN\_x4plus\_anime\_6B.pth} is selected instead, as its lighter 6-block architecture is better tuned to sharp-edged line patterns.

\subsection{DStretch Colour Decorrelation}

\textbf{\texttt{dstretch(img, colour\_space="LAB")}:} Implements the DStretch algorithm \cite{Harman2008} in pure Python using NumPy. The steps are:
\begin{enumerate}
    \item Convert the image to the specified colour space (LAB, YCbCr, or a Harman custom space such as YBK).
    \item Reshape the three-channel image to an $(N \times 3)$ pixel matrix $P$, where $N$ is the total pixel count.
    \item Compute the $3 \times 3$ covariance matrix $\Sigma = \frac{1}{N} P^T P - \bar{p}\bar{p}^T$, where $\bar{p}$ is the channel mean vector.
    \item Perform eigendecomposition: $\Sigma = Q \Lambda Q^T$, where $Q$ is the eigenvector matrix and $\Lambda = \text{diag}(\lambda_1, \lambda_2, \lambda_3)$.
    \item Whiten: $P_{\text{white}} = (P - \bar{p}) Q \Lambda^{-1/2}$.
    \item Re-correlate: $P_{\text{out}} = P_{\text{white}} \Lambda^{1/2} Q^T + \bar{p}$.
    \item Clip to $[0, 255]$ and convert back to BGR.
\end{enumerate}

The effect is to decorrelate the colour channels, amplifying colour differences invisible in the original image. This reveals inscription pigment on stone and rock surfaces where the colour contrast is below the threshold of human perception.

\subsection{Sharpening and Chain Orchestration}

\textbf{\texttt{sharpen(img, amount=1.5)}:} Computes a Gaussian-blurred version of the image (\texttt{sigma=1.0}), then adds the weighted difference to the original: $I_\text{out} = I + \text{amount} \cdot (I - \text{blur}(I))$. This unsharp mask recovers edge definition after denoising and super-resolution.

\textbf{\texttt{enhance(img\_path, output\_path, use\_dstretch=False)}:} The main entry point selects the per-artefact chain. Stone inscriptions and paper manuscripts run denoise $\rightarrow$ Real-ESRGAN $\rightarrow$ sharpen. Cave paintings run denoise $\rightarrow$ DStretch $\rightarrow$ sharpen. Palm leaf manuscripts run white balance $\rightarrow$ denoise $\rightarrow$ Real-ESRGAN.

\section{Stage 3 --- Binarisation (\texttt{src/binarise.py})}

\textbf{\texttt{binarise\_sauvola(img, window\_size=25)}:} Converts the image to grayscale, then calls scikit-image's \texttt{threshold\_sauvola()} with the specified window size. A binary mask is produced by comparing each pixel to its local Sauvola threshold. The binary image is inverted so that text pixels are dark and background is white, matching the standard convention for downstream processing.

\textbf{\texttt{binarise\_otsu(img)}:} Applies scikit-image's \texttt{threshold\_otsu()} to compute a global threshold from the full image histogram. This method is dispatched for paper manuscripts where the bimodal histogram assumption is valid.

\textbf{\texttt{binarise\_adaptive(img)}:} Calls \texttt{cv2.adaptiveThreshold} with the \texttt{cv2.ADAPTIVE\_THRESH\_MEAN\_C} mode and a block size of 11. This fallback is used for mixed-quality inputs where Sauvola produces artefacts due to extreme local contrast variation.

\textbf{\texttt{remove\_noise\_blobs(binary, min\_size=50)}:} Uses scikit-image's \texttt{remove\_small\_objects} to eliminate connected components with fewer than 50 pixels. Components this small correspond to dust, surface pitting, and scanning artefacts rather than character strokes.

\textbf{Morphological closing:} After noise removal, \texttt{cv2.morphologyEx} with \texttt{cv2.MORPH\_CLOSE} and a $2 \times 2$ rectangular kernel reconnects broken character strokes. This post-processing step is essential for weathered stone inscriptions where centuries of erosion have separated strokes that were once continuous.

\textbf{\texttt{binarise(img\_path, output\_path, method="sauvola")}:} The main entry point dispatches to the appropriate thresholding function, applies noise removal and morphological closing, and saves the result as a lossless PNG file.

\section{ECE Signal Processing and Quality Modules}

The ECE team contributes three source modules that extend the core CS/IT pipeline with signal-processing rigour. These modules are called from the pipeline runner as optional or automatic steps and are fully integrated into the Phase~1 system.

\subsection{Custom Filter Design (\texttt{src/filters.py})}

\textbf{\texttt{gabor\_filter\_bank(img, frequencies=[0.1, 0.2, 0.4], orientations=8)}:} Constructs a bank of $3 \times 8 = 24$ Gabor filters covering three spatial frequencies and eight equally spaced orientations ($0°$ to $157.5°$) \cite{Daugman1985}. Each filter is a complex-valued 2D Gabor kernel with Gaussian envelope. The convolution response is computed for every frequency--orientation pair, producing a 24-channel response tensor. Inscriptions exhibit strong directional edges (vertical and horizontal strokes in Devanagari, diagonal strokes in Tamil); the Gabor bank isolates character-background separation from background texture based on spatial frequency content alone, independent of intensity. The output is used to refine binarisation boundaries on difficult palm leaf and stone samples.

\textbf{\texttt{directional\_edge\_enhance(img, angle\_deg=45.0)}:} Applies a Gabor filter tuned to the specified orientation angle to enhance carving or stroke edges running in a given direction. For palm leaf manuscripts, where horizontal leaf-fibre texture interferes with vertical character strokes, a vertical-orientation Gabor pass sharpens the character edges while suppressing the horizontal fibre pattern.

\textbf{\texttt{remove\_periodic\_noise\_fft(img, threshold=0.1)}:} Converts the image to the frequency domain using the Fast Fourier Transform, identifies spectral peaks above the threshold (which correspond to periodic scanner line artefacts), zeros them using a notch filter mask, and applies the inverse FFT. Because text content is broadband and scanner artefacts are narrowband, this removes scanner stripes without affecting character information.

\subsection{Colour Channel and Histogram Analysis (\texttt{src/analysis.py})}

\textbf{Per-channel statistics:} Functions compute the mean, standard deviation, skewness, and kurtosis for each colour channel (R, G, B and L, a, b in LAB space) across the image. These statistics characterise the colour distribution of each artefact type and inform DStretch colour space parameter selection.

\textbf{DStretch colour space tuning:} Empirical analysis across 20 samples per artefact category determined optimal DStretch colour spaces: LAB and YDS for stone inscriptions, YBK for cave and rock art, and standard LAB for copper plates and paper. These per-material defaults are hard-coded in the \texttt{enhance.py} routing logic.

\textbf{Before/after histogram visualisation:} The \texttt{plot\_histogram\_comparison()} function generates side-by-side per-channel histogram plots of the preprocessed input and the enhanced output, used as report figures to demonstrate quantitative contrast improvement.

\subsection{Image Quality Metrics (\texttt{src/metrics.py})}

The metrics module is automatically called by the pipeline runner after Stage~2 enhancement. All four metrics are computed and logged for every processed image.

\textbf{\texttt{compute\_psnr(original, enhanced)}:} Wraps \texttt{skimage.metrics.peak\_signal\_noise\_ratio} to compute PSNR in decibels between the enhanced image and the original scan. Returns a single float.

\textbf{\texttt{compute\_ssim(original, enhanced)}:} Wraps \texttt{skimage.metrics.structural\_similarity} with \texttt{multichannel=True} to compute SSIM across all three colour channels and return the mean value.

\textbf{\texttt{compute\_cnr(img, text\_mask)}:} Accepts the enhanced image and a binary text mask (produced by Sauvola binarisation as a proxy for text pixel locations). Computes $\text{CNR} = |\mu_\text{text} - \mu_\text{bg}| / \sigma_\text{bg}$ where $\mu$ and $\sigma$ are computed from the text and non-text regions of the grayscale image.

\textbf{\texttt{compute\_sharpness(img)}:} Converts the image to grayscale, applies the Laplacian operator (\texttt{cv2.Laplacian}), and returns the variance of the Laplacian response. Higher variance indicates sharper edges.

\textbf{\texttt{full\_quality\_report(original, enhanced, text\_mask)}:} Convenience wrapper that calls all four metric functions and returns a dictionary \texttt{\{"psnr": ..., "ssim": ..., "cnr": ..., "sharpness": ...\}} for insertion into the processing log.

\section{Testing and Validation}

\subsection{Unit Tests}

\textbf{\texttt{tests/test\_preprocess.py}} contains four unit tests verifying: (1) CLAHE produces a histogram with a wider range than the original; (2) white balance reduces the maximum per-channel mean deviation; (3) border cropping removes rows and columns below the threshold; and (4) EXIF orientation is resolved so the output image dimensions match the visually correct orientation.

\subsection{Integration Tests}

\texttt{tests/test\_enhance.py} processes a known test image through the full Stage~1--2 chain and asserts that the PSNR of the enhanced output exceeds 30\,dB relative to the original.

\subsection{Continuous Integration}

A GitHub Actions workflow runs on every push to the \texttt{develop} and \texttt{main} branches. The workflow executes Python type checking with \texttt{mypy} across all source modules, then runs the full pytest suite (5 tests). Merge to \texttt{main} is blocked if any test fails or if type checking raises an error. This prevents stage regressions from propagating to the stable branch.

\section{Summary}

The Phase~1 implementation translates the Chapter~3 design into a tested, version-controlled system covering all five artefact categories across three pipeline stages. The preprocessing module establishes a reliable quality baseline through CLAHE, white balance, and EXIF correction. The enhancement module applies Real-ESRGAN super-resolution and DStretch colour decorrelation along artefact-specific chains. The binarisation module provides Sauvola adaptive thresholding as the primary method with morphological refinement. The ECE signal processing modules --- Gabor filter bank, FFT-based noise removal, colour analysis, and objective quality metrics --- integrate directly into the pipeline, ensuring every processed image is automatically scored. A 5-test continuous integration suite validates all components. Chapter~5 now evaluates the quality and performance of these outputs through objective image quality metrics.
